% \subsection{Introduction}
% \label{emt:sec:intro}

\begin{quote}
	``It so happens that a tree format is the only sane format...''~\cite{linux:torvalds_pt}
	---Linus Torvalds, 2002
\end{quote}

% \vspace{-5pt}

Virtual memory translation has become a major performance bottleneck
	of emerging memory-intensive computing~\cite{basu2013efficient,gupta2021midgard,
		skarlatos2020ecpt,park2022fpt,kanellopoulos2023utopia,margaritov2019asap,kwon2016coordinated}.
With unprecedented growth of memory capacity,
	driven by terabyte-scale memory~\cite{news:aws:18tb,news:gcloud:12tb} 
	and
	memory expanders like CXL~\cite{li2023pond,sharma2023cxl},
TLBs cannot scale in the same way as memory.
Moreover, emerging workloads like 
	machine learning, graph processing, and bioinformatics
	have irregular memory access patterns with weak locality, making
	TLBs and other MMU caches less efficient.
As a result, TLB misses are inevitably increasing, resulting in
	expensive address translation
	across the memory system.

However, today's translation architectures were 
	designed at a time
	of scarce memory, and optimize 
	space-efficiency over performance.
The {\it de facto} schemes organize translations
	into a multi-level radix tree~\cite{intel:man,amd:man,aarch64:man};
upon a TLB miss, the MMU must {\it sequentially} walk the tree, resulting in multiple memory accesses.
The x86-64 architecture uses a four-level
	tree, with a fifth level added in recent hardware~\cite{intel:5level,intel:5level:implement}. % like Intel Sunny Cove
In virtualized environments, translation overhead is magnified by nested translation~\cite{bhargava2008accelerating,merrifield2016performance}
	which takes
%	where the MMU performs 
	a two-dimensional 
	walk over the guest and host page tables, resulting in up to
% A nested translation takes up to 
24 sequential memory accesses on four-level page tables. % (35 with five levels).
%	and up to 35 with five-level page tables.


To address this pressing problem,
	many new hardware architectures for MMUs have been developed to realize fast translation
	for today's terabyte-scale, heterogeneous memory.
% ~\cite{skarlatos2020ecpt,gupta2021midgard,
%	park2022fpt,self:dmt,basu2013efficient,kanellopoulos2023utopia}.
For example, hashing-based translation schemes are 
	revisited~\cite{gosakan2023mosaic,jovan2022nestedecpt,skarlatos2020ecpt,jovan2023memory,yaniv2016hash,vmware2021hashing,Kwon24},
	as hashing is inherently more scalable than walking a tree~\cite{yaniv2016hash};
a recent hashing scheme, based on Elastic Cuckoo Page Table (ECPT), 
	is reported to reduce translation overhead significantly 
	by enabling parallel lookups of page table entries~\cite{skarlatos2020ecpt,jovan2022nestedecpt}.
%	and supports modern paging features like huge pages and page sharing.
New translation schemes using flattened or linear page tables~\cite{park2022fpt,karakostas2015redundant} have also been developed.
In addition, recent studies advocate for hybrid translation architectures
	that use different schemes collectively or selectively~\cite{alam2017doityourself,ahn2012revisiting,hoang2010nested,ghose2019pim,kanellopoulos2023utopia,self:dmt}.
%	and use them selectively according to applications~\cite{alam2017doityourself} and configurations~\cite{ahn2012revisiting,hoang2010nested,ghose2019pim}.

% Proposals for hybrid design for virtualization (e.g., tree for guest and hash table for host)
% or heterogeneous memory (e.g., tree for CXL).

Unfortunately, OS support for hardware innovations falls short;
	few aforementioned new hardware schemes were
	experimented with commodity OSes like Linux. % \TODO{and FreeBSD}.
Instead, evaluations of experimental architectures mostly use
	performance models to estimate OS overhead~\cite{self:dmt,margaritov2019asap,yaniv2016hash,karakostas2015redundant,alverti2020enhancing,ryoo2017rethinking},
	or trace-driven simulation with traces 
	collected by running workloads on unmodified Linux~\cite{skarlatos2020ecpt,ryoo2017rethinking,yaniv2016hash,margaritov2019asap}.
	% \jztodo{@siyuan: can you confirm on this?}.
The assumption is that OS overhead 
	on different translation architectures is constant.
However, our paper shows that translation architectures
	could have significant impacts on OS performance.

% Currently, most hardware designs were evaluated with trace-based simulation,
%	which hardly accounts for \TODO{[need strong and fair arguments.}
%	also need to account for full-system simulation]}
% even full-system simulation is down with a trace-driven approach;
In fact, the difficulties of OS support has affected hardware research---disruptive 
	hardware designs are often considered ``{\it undesirable}'' and lose to incremental approaches (see~\cite{margaritov2019asap}).
Our discussions with hardware vendors 
	tell us that the lack of commodity OS support and evaluation is a major barrier
	to assessing and adopting new hardware translation schemes.

% it is a time that is hard for clean-slate design that can ignore OS support.

The unsatisfactory OS effort is largely due to memory management systems in 
	commodity OSes not having the extensibility
	for emerging translation architectures.
For example, Linux 
    assumes a radix-tree based page table structure
    and lacks extensibility to support hardware schemes that cannot fit in its tree 
    definition.
As a result, supporting a new architecture often requires heavy
    modifications of memory management code
    in architecture-independent modules.
% This is a daunting task, given the size and complexity of modern memory management subsystems 
%	and their tight coupling with memory translation 
%	(\S\ref{emt:sec:background}).
Essentially, Linux provides no extensible
	interface for memory translation,
	unlike its
	other subsystems (e.g., VFS for file systems~\cite{linux:vfs}). 
	% (which enables diverse filesystems for different storage devices and usages).
% Despite the proposal of reusing VFS for memory management~\cite{tabatabai:hotos23},
%	VSF cannot provide native support for memory translation (\S\ref{emt:sec:need}).
% is caused by 
%	historical views of memory subsystem, as reflected by the quotes from 
%	Torvalds in 2002.
% However, emerging memory systems have drastically different characteristics
%	and we argue that modern OSes must provide extensibility to accommodate
%	terabyte-scale and heterogeneous memory.
% BSD kernels~\cite{cranor99uvm} %that adopt Mach's design~\cite{rashid1987machine,rashid1987machine}
%	have more extensible memory management systems (from Mach's design~\cite{rashid1987machine});
% but, their address map interface is
%	not expressive enough for hardware-specific optimizations
%	in OS memory management (\S\ref{emt:sec:need}).

% Modern mm subsystem is very complex with many features (e.g., THP, PTI)
%	and their implementations are tightly coupled with address translation.
% Few of them were considered by architecture research
%	but can have profound implications of practical adoptions.

% An important roadblock is that there's no generic VM design to implement them in the software stack.
% The reason is there is no generic VM design to implement and evaluate them.

\para{Contributions.}
We develop a pragmatic OS framework and toolchains
	atop Linux to embrace 
	new hardware translation architectures for today's memory technologies.
We term our framework 
%	such as x86-like radix tree and hash table.
	Extensible Memory Translation (EMT).
We target Linux %, instead of OSes that are easier to extend (e.g., microkernels), 
	as it is 
	still the {\it de facto} commodity OS and 
	is mostly assumed by architecture research on memory systems.

EMT provides an architecture-neutral Linux interface that 
	1) supports diverse memory translation architectures,
    2) enables hardware-specific optimizations,
	3) accommodates modern hardware and OS complexity,
	and 4) has negligible overhead over Linux's hardwired implementation.
% \schaix{With EMT, new translation architectures can be effectively supported
% 	and experimented on Linux.}
\schai{
By abstracting translation-related operations, 
	EMT simplifies the effort to support and experiment with new translation architectures on Linux.
% \weiwei{I think "EMT enables easier..." can be removed as it says the same thing as the previous one.}
% EMT enables easier support for new translation architectures on Linux.
% EMT avoids developers from refactoring memory management code in Linux "architecture-independen" code.
With EMT, developers only need to implement an MMU driver
	as a hardware-specific implementation of translation logic,
	without modifying the architecture-independent memory management code.
% supporting a new memory-translation architecture only requires developers to write a MMU driver, 
% 	a architecture-dependent implementation of memory transltion logic. 
% EMT avoids changing the architecture-independent memory management code. 
}
EMT also makes it easy to profile and analyze OS performance with regard to 
	hardware translation schemes as it abstracts translation-related
	% memory management
	operations. %  related to translation.
% 	by adding hardware-specific code,
%	while reusing architecture independent modules.
% EMT is inspired by the pmap interface in BSD/Mach
%	for separating machine dependent and independent code~\cite{rashid1987machine,rashid1987machine},
%	but is different.
% Different from pmap as a general representation of physical address maps, 
%	EMT targets on
%	enabling OS support for new translation architectures:
%	it exposes low-level primitives to OS memory management %like pointers to translation entries
%	and enables hardware-specific customizations on architecture-neutral implemenations. % on primitives iterators and locks.
% Built on Linux,
%	EMT minimizes 
%	indirection overhead and avoids introducing new abstractions.

% \TODO{I feel I miss something; the writing is not as exciting as I hope to be.
% We need to discuss and brainstorm.}

We implement EMT on Linux (v5.15), referred to EMT-Linux. 
% We port Linux onto EMT,
%	referred to as EMT-Linux.
We modularize
	architecture-independent code with EMT API, removing 
	hardwired assumptions.
EMT realizes negligible overhead % \schai{(on average less than 0.5\%)}
	through careful engineering and optimizations
	\schai{(with less than 0.5\% overhead 
	on average across key OS operations)}.
%	with cache efficiency.
Moreover, EMT-Linux realizes all existing features and 
	hardware-specific optimizations
	in vanilla Linux.
% We then refactor hardware-specific memory management code using EMT.
% \TODO{???: by in-place adaptation without runtime overhead.}
	
% on both micro and macro benchmarks.

% \jyk{Adding number.}
% ---an extensible interface can be accomplished without sacrificing system performance.
% We used Linux Testing Project (LTP) to check the correctness of EMT-Linux.

We build on EMT-Linux to add OS support for ECPT~\cite{skarlatos2020ecpt} and FPT~\cite{park2022fpt}, 
	two new translation schemes.
%	hashing-based translation scheme 
%	different from x86-64's scheme.
Porting these new hardware schemes on Linux without a framework like EMT would require 
	major rewriting of Linux's
	memory management module.
	% including architecture-independent portion which is hardwired to tree-based translation. 
With EMT, supporting them on Linux is modularized 
% by implementing basic functions and customized optimizations 
	with manageable engineering efforts.

Evaluating new hardware schemes on EMT-Linux faces a common challenge 
	of hardware-software codesign---the hardware is not yet available
	to run the OS.
To address this problem, we assemble a toolchain that 
	runs EMT-Linux on QEMU with an emulated MMU where we implement
	the hardware translation logic.
Our toolchain enables us to understand the OS perspective of new translation
	architectures,
	such as its OS overhead over x86. %  as well as other translation . 
The toolchain also supports cycle-accurate hardware simulation.
%   using existing simulators.
%	recording memory traces of EMT-Linux
%	and replaying them in hardware simulators.

We share our experience of supporting new translation schemes
	on EMT-Linux, which
	enables us to understand OS memory management challenges 
	% of hardware translation architecture,
	% to effectively manage the new, fast translation architecture
	beyond hardware perspectives.
We present our reflection on the ECPT design and address correctness challenges such as managing
	the kernel page table (kECPT) and the paradox involved in 
	changing the translations of the kECPT itself and of the kernel code 
	that manages kECPT, as well as performance challenges like 
	efficient locking and memory scanning.
%	but have to be addressed for practical use cases.
% elegant hardware designs may not transfer to software implementations.
Arguably, OS framework support is essential to encourage and embrace 
	disruptive hardware innovations,
	and EMT is an important step forward.
%	we find current address translation designs overlook design requirements from the OS perspectives.
% Also, we find current designs often overlook design aspects, as they are buried deep in the code mountain.

% It's time.	
% Statements that only tree is the practical solution;

\para{Summary.} This paper makes the following contributions:
\begin{packed_itemize}%[\labelitemi]
	\item A discussion on empowering new, experimental hardware-assisted translation architectures
		on commodity OSes.
	\item EMT as an extensible framework for developing and evaluating OS memory management 
		on new translation architectures,
		and its implementation on Linux.
	\item An experience of building ECPT and FPT on Linux using EMT and the reflection 
		on hardware/OS designs.
	\item An open platform for developing, testing, and evaluating OS kernels 
		on new memory translation architectures.
	% \item \schaix{All code and data will be released, including EMT-Linux, ECPT/FPT
	% implementation, and the emulator toolchains.}
	\item \schai{Our artifacts: \url{https://github.com/xlab-uiuc/emt}.}
\end{packed_itemize}



% While the marvelous hardware designs strive to improve performance of address translation,
%    few of them consider the challenges of integrating the designs into a modern operating system. 
% To help future researchers and engineers understand the design space in OS,
%    we delve into Linux to study its structure and design, 
%    memory and page table management in particular.